% Provenance (never printed): see sources/outline_sources.md (Year 2: radiative forcing
% quantified; feedbacks and climate sensitivity; detection and attribution; scenarios; the
% remaining carbon budget; components of sea-level rise; physical science only). Every statement
% from IPCC AR6 WG1 (SPM; chapters 3, 5, 7 and 9; Annex III data), ledger rows in
% sources/ledger/book2.md (ipcc:*) and sources/ledger/book4.md (ar6*:*). Curves and numbers from
% figdata/bachelor-2/anthropogenic.py and the Book 4 second-writer scratch calculator calc/ch27.py.
\chapter{Perubahan Iklim Antropogenik}\label{ch:b2:anthropogenic-climate-change}

Sejak 1970 permukaan global telah menghangat lebih cepat daripada dalam periode lima puluh tahun lain mana pun selama
sedikitnya dua ribu tahun terakhir (\emph{keyakinan tinggi}). Volume-volume sekolah telah menceritakan apa yang berubah dan
mengapa. Bab ini mengubah cerita itu menjadi suatu anggaran: berapa watt per meter persegi yang ditambahkan oleh
pendorong-pendorong manusia, berapa banyak yang dikirim balik ke angkasa oleh planet yang lebih hangat, berapa pemanasan yang dihasilkan oleh setiap ton karbon
dioksida, dan bagaimana kenaikan muka laut terbagi antara air yang menghangat dan es yang mencair. Setiap angka
disertai ketidakpastiannya, dalam bahasa terkalibrasi dari penilaian-penilaian ilmiah.
% ledger: ipcc:fastest-2000yr, ipcc:likelihood

\begin{recall}
\textbf{Dari volume-volume sekolah}: gaya radiatif, $\qty{3.93}{W.m^{-2}}$ untuk satu penggandaan
karbon dioksida, dan hukum logaritmiknya; umpan balik iklim; \omterm{def:b2:anthropogenic-climate-change:ecs}{sensitivitas iklim kesetimbangan}, dengan perkiraan
terbaik \qty{3}{\celsius}, \emph{mungkin} 2.5 sampai \qty{4}{\celsius}; \omterm{def:b2:anthropogenic-climate-change:da}{atribusi}, skenario emisi,
dan proyeksi; bahasa terkalibrasi: \emph{mungkin} berarti 66 sampai \qty{100}{\percent}, \emph{sangat
mungkin} 90 sampai \qty{100}{\percent}, dan suatu rentang $[x \text{ sampai } y]$ adalah rentang yang sangat mungkin. Pengaruh
manusia telah menghangatkan iklim sebesar \qty{1.07}{\celsius} antara 1850--1900 dan 2010--2019.
\textbf{Dari buku-buku fisika}: benda hitam dan hukum Stefan--Boltzmann; pemuaian termal.
\textbf{Dari buku-buku biologi}: siklus karbon, fraksi yang tetap di udara, penyerap karbon, dan
anggaran karbon biosfer. \textbf{Dari \cref{ch:b2:deep-time-climate}}: siklus karbon lambat,
yang mengambil kembali kelebihan karbon dioksida hanya selama ratusan ribu tahun.
\end{recall}
% ledger: ipcc:erf-2xco2, ipcc:log-co2, ipcc:ecs, ipcc:likelihood, ipcc:attribution

\section{Gaya radiatif secara kuantitatif}\label{sec:b2:anthropogenic-climate-change:forcing}

\begin{definition}[Gaya radiatif efektif]\label{def:b2:anthropogenic-climate-change:erf}
Yang disebut \emph{gaya radiatif efektif}\index{gaya radiatif efektif} (ERF) dari suatu gangguan
--- perubahan suatu gas rumah kaca, aerosol, atau Matahari --- adalah perubahan yang ditimbulkannya pada fluks bersih radiasi
ke bawah di puncak atmosfer, dalam \unit{W.m^{-2}}, setelah atmosfer,
awan, dan permukaan daratan menyesuaikan diri terhadapnya, tetapi sebelum suhu permukaan global
menanggapinya. Penyesuaian yang dihitung adalah penyesuaian yang disebabkan oleh agen itu sendiri, tidak bergantung pada
suhu permukaan.
\end{definition}
% ledger: ar6ch7:erf-def

Definisi itu menarik garis antara gaya radiatif dan tanggapan berdasarkan apa yang bergantung pada suhu
permukaan, bukan berdasarkan kecepatan. Gaya radiatif seketika dari karbon dioksida adalah perubahan fluks tepat
setelah gas itu ditambahkan; stratosfer, yang kemudian mendingin dalam beberapa bulan, mengubahnya untuk pertama kali;
troposfer, uap airnya, dan awannya mengubahnya lagi. Untuk penggandaan karbon
dioksida penilaian itu mendapatkan \qty{3.75}{W.m^{-2}} setelah stratosfer menyesuaikan diri, dan
penyesuaian di troposfer yang sebagian saling meniadakan: \qty{-0.60}{W.m^{-2}} dari pemanasannya, +0.22
dari uap air, +0.45 dari awan, dan +0.11 dari permukaan daratan, \qty{+0.18}{W.m^{-2}} seluruhnya.
ERF-nya adalah $3.75 + 0.18 = \qty{3.93}{W.m^{-2}}$, \qty{6}{\percent} lebih besar daripada
\qty{3.71}{W.m^{-2}} dalam penilaian sebelumnya (\cref{exo:b2:anthropogenic-climate-change:1}).
% ledger: ar6ch7:2xco2, ar6ch7:erf-def

Gaya radiatif setiap agen sejak 1750 telah direkonstruksi tahun demi tahun. Gambar di bawah menunjukkannya:
karbon dioksida pada \qty{2.16}{W.m^{-2}} pada 2019, gas-gas rumah kaca lainnya --- metana,
dinitrogen oksida, gas-gas terhalogenasi, dan ozon --- pada \qty{1.69}{W.m^{-2}} bersama-sama, dan
aerosol pada \qty{-1.06}{W.m^{-2}}. Gaya radiatif manusia total adalah \qty{2.72}{W.m^{-2}}, dan rentang
sangat mungkinnya, 1.96 sampai \qty{3.48}{W.m^{-2}}, lebar terutama karena aerosol. ERF aerosol selama
1750--2014 dinilai sebesar $-1.3$ [$-2.0$ sampai $-0.6$] \unit{W.m^{-2}} (\emph{keyakinan sedang}),
sebagian besar melalui pengaruhnya pada awan, yang menjadi lebih cerah bila aerosol menyediakan lebih banyak butir air.
Letusan gunung api menambahkan lonjakan negatif yang singkat (\cref{ch:b2:eruptive-styles}): yang terdalam dalam
rekaman itu, \qty{-4.8}{W.m^{-2}} pada 1816, menyusul letusan Tambora pada 1815.
% ledger: data:b4-ar6-erf, ipcc:erf-total, ipcc:erf-agents, ar6ch7:aerosol, usgs:msh-comparisons

\begin{omfigure}
\centering
\begin{tikzpicture}
\begin{axis}[omrecord, width=0.92\linewidth, height=6.2cm, xmin=1750, xmax=2020, ymin=-5.2,
  ymax=3.8, xlabel={tahun}, ylabel={ERF sejak 1750 (\unit{W.m^{-2}})},
  xtick={1750,1800,...,2000}, /pgf/number format/1000 sep={}, xtick pos=bottom, ytick pos=left,
  tick label style={font=\scriptsize}, label style={font=\small},
  legend style={at={(0.5,-0.24)}, anchor=north, font=\scriptsize, draw=black!30,
  legend columns=3, /tikz/every even column/.append style={column sep=6pt}},
  legend cell align=left]
% computed in figdata/bachelor-2/anthropogenic.py part erf; ledger: data:b4-ar6-erf
\addplot[name path=hi, draw=none, forget plot] table[x=year, y=hi] {figdata/out/bachelor-2-anthropogenic-erf.dat};
\addplot[name path=lo, draw=none, forget plot] table[x=year, y=lo] {figdata/out/bachelor-2-anthropogenic-erf.dat};
\addplot[black!15, forget plot] fill between[of=hi and lo];
\addplot[black!45, line width=0.4pt] table[x=year, y=natural] {figdata/out/bachelor-2-anthropogenic-erf.dat};
\addlegendentry{alami (gunung api, Matahari)}
\addplot[omThm, line width=1pt] table[x=year, y=co2] {figdata/out/bachelor-2-anthropogenic-erf.dat};
\addlegendentry{karbon dioksida}
\addplot[omThm!50!orange, line width=1pt, densely dashed] table[x=year, y=other_ghg] {figdata/out/bachelor-2-anthropogenic-erf.dat};
\addlegendentry{gas rumah kaca lain dan ozon}
\addplot[omDef, line width=1pt] table[x=year, y=aerosol] {figdata/out/bachelor-2-anthropogenic-erf.dat};
\addlegendentry{aerosol}
\addplot[black, line width=1.4pt] table[x=year, y=anthro] {figdata/out/bachelor-2-anthropogenic-erf.dat};
\addlegendentry{total manusia (abu-abu: 5--95\,\%)}
\draw[black!40] (axis cs:1750,0) -- (axis cs:2020,0);
\node[anchor=west, font=\scriptsize, text=black!60, inner sep=1pt] at (axis cs:1817.5,-4.75) {$\leftarrow$ Tambora};
\end{axis}
\end{tikzpicture}
\par\smallskip{\small \omterm{def:b2:anthropogenic-climate-change:erf}{Gaya radiatif efektif} dari 1750 sampai 2019, relatif terhadap 1750, menurut pendorongnya.
Gaya radiatif manusia total mencapai \qty{2.72}{W.m^{-2}} pada 2019, dengan rentang sangat mungkinnya diarsir.
Data: Laporan Penilaian Keenam IPCC, Kelompok Kerja~I, data tambahan Lampiran~III (CC BY 4.0).}
\end{omfigure}
% values: anthro 2.7204 (1.9569-3.4819), co2 2.1563, other_ghg 1.6850, aerosol -1.0583, volcanic -4.8162 in 1816 (anthropogenic.py numbers)

\section{Anggaran energi}\label{sec:b2:anthropogenic-climate-change:energy}

Gaya radiatif yang positif membuat planet memperoleh energi. Ketika menghangat, planet itu memancarkan lebih banyak ke angkasa, dan
perolehannya menyusut; iklim baru menetap ketika keduanya seimbang.

\begin{definition}[Ketidakseimbangan energi Bumi; parameter umpan balik iklim]\label{def:b2:anthropogenic-climate-change:eei}
Yang disebut \emph{ketidakseimbangan energi Bumi}\index{ketidakseimbangan energi Bumi} $N$ adalah fluks bersih energi yang masuk ke
sistem iklim di puncak atmosfer, dirata-ratakan atas seluruh bola Bumi dan selama satu tahun atau lebih,
dalam \unit{W.m^{-2}}: sinar Matahari yang diserap dikurangi inframerah yang keluar. Besaran ini adalah laju perolehan kalor oleh samudra,
daratan, es, dan udara, per meter persegi permukaan Bumi.
\end{definition}

\begin{definition}[Parameter umpan balik iklim]\label{def:b2:anthropogenic-climate-change:lambda}
Yang disebut \emph{parameter umpan balik iklim}\index{parameter umpan balik iklim} $\lambda$ adalah perubahan
fluks bersih ke bawah di puncak atmosfer per derajat pemanasan permukaan global, dalam
\unit{W.m^{-2}.\celsius^{-1}}, yang disebabkan oleh tanggapan sistem iklim terhadap pemanasan itu. Iklim
yang stabil memiliki $\lambda < 0$: pemanasan menambah kehilangan energi ke angkasa.
\end{definition}

Volume sekolah menuliskan keseimbangan itu dengan besarnya $\lambda$, suatu bilangan positif; mulai dari sini
tandanya dipertahankan, seperti dalam penilaian-penilaian itu, yang menuliskan parameter ini sebagai $\alpha$.
% ledger: ar6ch7:energy-budget

\begin{proposition}[Keseimbangan energi global]\label{prop:b2:anthropogenic-climate-change:balance}
Misalkan $\Delta F$ \omterm{def:b2:anthropogenic-climate-change:erf}{gaya radiatif efektif} dan $\Delta T$ pemanasan permukaan global
relatif terhadap suatu keadaan setimbang. Jika tanggapan fluks keluar linear terhadap $\Delta T$,
\[
  \Delta N = \Delta F + \lambda\,\Delta T .
\]
Pada kesetimbangan yang baru $\Delta N = 0$, sehingga $\Delta T_{\text{eq}} = -\Delta F/\lambda$.
\end{proposition}
\begin{proof}
Menurut definisi, gaya radiatif mengubah fluks bersih ke bawah sebesar $\Delta F$ pada suhu permukaan
tetap, dan tanggapan terhadap pemanasan mengubahnya sebesar $\lambda\,\Delta T$ sampai orde pertama;
ketidakseimbangannya adalah jumlah keduanya. Dengan menyamakannya dengan nol diperoleh kesetimbangan.
\end{proof}

Mengukur $N$ berarti mengukur kalor yang diperoleh sistem, dan samudra menampung hampir seluruhnya.
Dari 1971 sampai 2018 sistem iklim memperoleh \qty{435}{ZJ} (zettajoule, \qty{e21}{J}): samudra
\qty{91}{\percent}, yaitu \qty{242}{ZJ} di \qty{700}{m} teratasnya dan \qty{154}{ZJ} di bawahnya; daratan
\qty{5}{\percent}, es yang mencair \qty{3}{\percent}, dan atmosfer \qty{1}{\percent}. Ketidakseimbangan
rata-ratanya adalah \qty{0.57}{W.m^{-2}} selama 1971--2018, dan nilai itu telah bertambah: \qty{0.79}{W.m^{-2}}
[0.52 sampai 1.06] selama 2006--2018 (\cref{exo:b2:anthropogenic-climate-change:2}).
% ledger: ar6ch7:inventory, ipcc:heating

\begin{omfigure}
\centering
\begin{tikzpicture}[font=\small, >=Stealth]
\draw[fill=omDef!8, draw=omDef!60] (-3.2,0) rectangle (3.2,1.6);
\node[align=center] at (0,1.05) {sistem iklim};
\node[align=center, font=\scriptsize] at (0,0.5) {samudra \qty{91}{\percent}, daratan \qty{5}{\percent}, es \qty{3}{\percent}, udara \qty{1}{\percent}};
\draw[->, line width=2.2pt, omThm] (-2.0,3.3) -- node[pos=0.12, left=3pt, align=right, font=\scriptsize] {gaya radiatif $\Delta F$} (-2.0,1.7);
\draw[->, line width=2.2pt, omDef] (2.0,1.7) -- node[pos=0.88, right=3pt, align=left, font=\scriptsize] {tanggapan $\lambda\,\Delta T$} (2.0,3.3);
\draw[->, line width=1.4pt, black!70] (0,-0.1) -- node[right=3pt, align=left, font=\scriptsize] {tersimpan: $\Delta N = \Delta F + \lambda\,\Delta T$} (0,-1.0);
\draw[black!40, dashed] (-3.4,2.4) -- (3.4,2.4) node[right, font=\scriptsize, text=black!60] {puncak atmosfer};
\end{tikzpicture}
\par\smallskip{\small Anggaran energi suatu iklim yang dipaksa (skematis). Gaya radiatif menambahkan energi di
puncak atmosfer, tanggapan planet yang lebih hangat membuang sebagian ($\lambda\,\Delta T$ bernilai
negatif), dan selisihnya tersimpan, sebagian besar di samudra. Bagian-bagian kalor yang diperoleh dari 1971
sampai 2018.}
\end{omfigure}
% ledger: ar6ch7:inventory

\begin{method}[Sensitivitas iklim dari anggaran energi historis]\label{met:b2:anthropogenic-climate-change:energy-budget}
\begin{enumerate}
  \item Pilihlah dua periode, satu acuan yang dekat dengan kesetimbangan dan satu yang baru; ambillah perubahan
        gaya radiatif $\Delta F$, ketidakseimbangan $\Delta N$, dan suhu global $\Delta T$.
  \item Selesaikan keseimbangan itu untuk parameter umpan balik: $\lambda = (\Delta N - \Delta F)/\Delta T$.
  \item Sensitivitas terhadap penggandaan kemudian adalah $-F_{2\times}/\lambda$; pemanasan pada saat
        gaya radiatif itu, yang diskalakan ke $F_{2\times}$, adalah $F_{2\times}\Delta T/\Delta F$.
\end{enumerate}
\end{method}

Dari 1850--1900 sampai 2006--2019 penilaian itu mengambil $\Delta N = \qty{0.59}{W.m^{-2}}$, $\Delta T =
\qty{1.03}{\celsius}$, dan $\Delta F = \qty{2.20}{W.m^{-2}}$. Maka $\lambda = (0.59 - 2.20)/1.03 =
\qty{-1.56}{W.m^{-2}.\celsius^{-1}}$, dan sensitivitasnya adalah $3.93/1.56 = \qty{2.5}{\celsius}$,
nilai yang didapatkan penilaian itu, dengan rentang sangat mungkin 1.6 sampai \qty{4.8}{\celsius}. Namun ini adalah
sensitivitas \emph{efektif}, untuk umpan balik pemanasan sejauh ini. Pemanasan paling
kuat di tempat umpan balik paling negatif, di Pasifik tropis bagian barat, dan lebih lemah di atas
Pasifik timur dan Samudra Selatan; seiring berkembangnya pola itu, parameter umpan balik diperkirakan
menjadi kurang negatif (\emph{keyakinan tinggi}), sehingga sensitivitas jangka panjang lebih tinggi daripada
perkiraan historis itu (\cref{exo:b2:anthropogenic-climate-change:7}).
% ledger: ar6ch7:energy-budget, ar6ch7:pattern, ipcc:erf-2xco2
% values: lambda_energy -1.563, ecs_energy 2.514 (anthropogenic.py numbers)

\section{Umpan balik dan sensitivitas}\label{sec:b2:anthropogenic-climate-change:feedbacks}

\begin{definition}[Respons Planck; umpan balik uap air, laju lapse, dan awan]\label{def:b2:anthropogenic-climate-change:feedbacks}
Parameter umpan balik dipecah menjadi sumbangan besaran-besaran yang berubah bersama
suhu permukaan. Yang disebut \emph{respons Planck}\index{respons Planck} adalah tambahan inframerah
yang dipancarkan ke angkasa ketika permukaan dan seluruh troposfer menghangat dengan besar yang sama. Yang disebut
\emph{umpan balik uap air}\index{umpan balik uap air} adalah perubahan fluks akibat bertambahnya
uap air atmosfer seiring suhu. Yang disebut \emph{umpan balik laju lapse}\index{umpan balik laju lapse}
adalah perubahan fluks akibat pemanasan yang tidak seragam menurut ketinggian. Yang disebut \emph{umpan
balik awan}\index{umpan balik awan} adalah perubahan fluks akibat perubahan jumlah, ketinggian, dan
kecerahan awan.
\end{definition}

\begin{proposition}[Respons Planck]\label{prop:b2:anthropogenic-climate-change:planck}
Jika Bumi memancarkan ke angkasa suatu fraksi tetap $1 - g$ dari inframerah $\sigma T^4$ permukaannya pada
suhu $T$, \omterm{def:b2:anthropogenic-climate-change:feedbacks}{respons Planck-nya} adalah
\[
  \lambda_{\text P} = -(1 - g)\,4\sigma T^3 .
\]
Dengan $g = 0.4$ dan $T = \qty{288}{K}$, $\lambda_{\text P} = \qty{-3.25}{W.m^{-2}.K^{-1}}$.
\end{proposition}
\begin{proof}
Fluks keluar adalah $(1 - g)\sigma T^4$; turunannya terhadap $T$ adalah $(1 - g)\,4\sigma
T^3$, dan kehilangan tambahan merupakan sumbangan negatif bagi fluks bersih ke bawah. Secara numerik $0.6
\times 4 \times \num{5.670e-8} \times 288^3 = 3.25$.
\end{proof}
% ledger: ar6ch7:planck, const:sigma, ipcc:budget-surface
% values: planck -3.2509 (anthropogenic.py numbers)

Fraksi $g$ adalah efek rumah kaca yang dinormalkan: atmosfer menahan \qty{159}{W.m^{-2}} dari
\qty{398}{W.m^{-2}} yang dipancarkan permukaan, $g \approx 0.4$. Penilaian itu, dengan kode radiasi
yang terperinci, mendapatkan $-3.22$ [$-3.4$ sampai $-3.0$] (\emph{keyakinan tinggi}). Uap air mengikuti
hukum Clausius--Clapeyron pada skala global, dan lebih banyak uap menyerap lebih banyak inframerah: \omterm{def:b2:anthropogenic-climate-change:feedbacks}{umpan
balik uap air} sangat positif, \qty{1.85}{W.m^{-2}.\celsius^{-1}} menurut satelit.
Di daerah tropis troposfer atas menghangat lebih banyak daripada permukaan, lalu memancarkan lebih banyak: \omterm{def:b2:anthropogenic-climate-change:feedbacks}{umpan
balik laju lapse} negatif di sana, positif di dekat kutub, dan sekitar $-0.50$ secara global. Keduanya
dinilai bersama, \num{1.30} [1.1 sampai 1.5], karena apa yang dimasukkan suatu model ke dalam yang satu cenderung
diambilnya dari yang lain. Salju dan es yang mencair menggelapkan permukaan, +0.35; awan menambahkan \num{0.42}
[$-0.10$ sampai $0.94$], ketidakpastian terbesar. \omterm{def:b2:anthropogenic-climate-change:feedbacks}{Umpan balik awan} bersih yang negatif \emph{sangat kecil kemungkinannya}.
% ledger: ar6ch7:planck, ar6ch7:wv-lr, ar6ch7:table710, ipcc:cloud-negative, ipcc:budget-surface

\begin{omfigure}
\centering
\begin{tikzpicture}
\begin{axis}[xbar, width=0.78\linewidth, height=5.4cm, xmin=-3.7, xmax=1.8, ymin=0.4, ymax=6.6,
  ytick={1,...,6}, yticklabels={bersih, biogeofisika dan kimia, awan, albedo permukaan,
  uap air + laju lapse, Planck}, y tick label style={font=\scriptsize},
  tick label style={font=\scriptsize}, xlabel={parameter umpan balik (\unit{W.m^{-2}.\celsius^{-1}})},
  label style={font=\small}, bar width=7pt, axis x line=bottom, axis y line=left,
  axis line style={-}, xmajorgrids, grid style={gray!22, line width=0.3pt}, extra x ticks={0},
  extra x tick style={grid style={black!60, line width=0.5pt}}]
% ledger: ar6ch7:table710 (central values and very likely ranges)
\addplot[fill=omDef!55, draw=omDef!80!black, error bars/.cd, x dir=both, x explicit,
  error bar style={black!70}] coordinates {
  (-3.22,6) += (0.22,0) -= (0.18,0)
  (-0.01,2) += (0.26,0) -= (0.26,0)
  (-1.16,1) += (0.65,0) -= (0.65,0)};
\addplot[fill=omThm!55, draw=omThm!80!black, error bars/.cd, x dir=both, x explicit,
  error bar style={black!70}] coordinates {
  (1.30,5) += (0.20,0) -= (0.20,0)
  (0.35,4) += (0.25,0) -= (0.25,0)
  (0.42,3) += (0.52,0) -= (0.52,0)};
\end{axis}
\end{tikzpicture}
\par\smallskip{\small Umpan balik iklim yang dinilai oleh IPCC (Laporan Penilaian Keenam, Kelompok
Kerja~I, Tabel 7.10), nilai tengah dan rentang sangat mungkin. Umpan balik positif memperkuat pemanasan
(merah), yang negatif meredamnya (biru); nilai bersih adalah jumlahnya.}
\end{omfigure}

\begin{definition}[Sensitivitas iklim kesetimbangan]\label{def:b2:anthropogenic-climate-change:ecs}
Yang disebut \emph{sensitivitas iklim kesetimbangan}\index{sensitivitas iklim kesetimbangan} (ECS) adalah
pemanasan permukaan global, pada kesetimbangan yang baru, yang disebabkan oleh penggandaan karbon dioksida
udara dari tingkat praindustrinya, dengan lembaran es dibuat tetap.
\end{definition}

\begin{proposition}[Sensitivitas dan umpan balik]\label{prop:b2:anthropogenic-climate-change:ecs}
Jika parameter umpan balik konstan, $\text{ECS} = -F_{2\times}/\lambda$. Dengan $F_{2\times} =
\qty{3.93}{W.m^{-2}}$ dan $\lambda = \qty{-1.16}{W.m^{-2}.\celsius^{-1}}$ bersih dari umpan balik
di atas, $\text{ECS} = \qty{3.4}{\celsius}$; \omterm{def:b2:anthropogenic-climate-change:feedbacks}{respons Planck} saja akan memberikan $3.93/3.22 =
\qty{1.2}{\celsius}$.
\end{proposition}
\begin{proof}
Terapkan \cref{prop:b2:anthropogenic-climate-change:balance} pada kesetimbangan dengan $\Delta F =
F_{2\times}$.
\end{proof}
% ledger: ar6ch7:table710, ipcc:erf-2xco2
% values: ecs_process 3.388 (anthropogenic.py numbers)

Umpan balik adalah satu jalur bukti; pemanasan sejak 1850, anggaran energi, dan iklim
masa lalu adalah jalur-jalur lainnya. Bila digabungkan, semuanya memberikan ECS hasil penilaian sebesar \qty{3}{\celsius}, \emph{mungkin}
2.5 sampai \qty{4}{\celsius} (\emph{keyakinan tinggi}), \emph{sangat mungkin} 2 sampai \qty{5}{\celsius}.
Karena ECS berubah sebagai $1/\lambda$, ketidakpastian yang simetris pada $\lambda$ memberikan rentang
ECS yang memanjang ke arah nilai tinggi (\cref{exo:b2:anthropogenic-climate-change:12}).
% ledger: ipcc:ecs

\begin{definition}[Respons iklim transien]\label{def:b2:anthropogenic-climate-change:tcr}
Yang disebut \emph{respons iklim transien}\index{respons iklim transien} (TCR) adalah pemanasan permukaan
global pada saat karbon dioksida udara telah berlipat dua, bila karbon dioksida itu naik \qty{1}{\percent}
per tahun dari tingkat praindustrinya, sekitar tujuh puluh tahun setelah permulaannya.
\end{definition}

TCR lebih kecil daripada ECS karena samudra masih menyerap kalor ketika penggandaan itu
tercapai. Perkiraan terbaiknya adalah \qty{1.8}{\celsius}, \emph{mungkin} 1.4 sampai \qty{2.2}{\celsius}
(\emph{keyakinan tinggi}); anggaran energi historis saja juga memberikan $3.93 \times 1.03/2.20 =
\qty{1.8}{\celsius}$.
% ledger: ar6ch7:tcr, ar6ch7:energy-budget
% values: tcr_energy 1.840 (anthropogenic.py numbers)

\section{Deteksi dan atribusi}\label{sec:b2:anthropogenic-climate-change:attribution}

\begin{definition}[Deteksi; atribusi; sidik jari iklim]\label{def:b2:anthropogenic-climate-change:da}
Yang disebut \emph{deteksi}\index{deteksi (perubahan iklim)} suatu perubahan iklim adalah menunjukkan bahwa
iklim telah berubah dalam suatu pengertian statistik yang terdefinisi, melampaui apa yang dapat dihasilkan oleh variabilitas internalnya,
tanpa menyebutkan penyebabnya. Yang disebut \emph{atribusi}\index{atribusi} adalah menilai
sumbangan beberapa penyebab terhadap suatu perubahan yang terdeteksi, dengan kemungkinan yang ditetapkan. Yang disebut \emph{sidik jari
iklim}\index{sidik jari iklim} adalah pola dalam ruang dan waktu dari tanggapan terhadap satu
pendorong, sebagaimana disimulasikan oleh model iklim, yang dengannya pendorong itu dikenali dalam pengamatan.
\end{definition}

Pendorong yang berbeda meninggalkan sidik jari yang berbeda. Tambahan gas rumah kaca menghangatkan troposfer dan mendinginkan
stratosfer, karena selimut inframerah yang lebih tebal membuat lapisan-lapisan atas memancarkan lebih banyak ke angkasa;
Matahari yang lebih terang akan menghangatkan keduanya. Pengamatan menunjukkan pola yang pertama: troposfer telah menghangat
sejak sedikitnya tahun 1950-an, dan \emph{hampir pasti} bahwa stratosfer telah mendingin. \emph{Sangat
mungkin} bahwa gas rumah kaca dari manusia adalah pendorong utama pemanasan troposfer
sejak 1979, dan \emph{amat sangat mungkin} bahwa penipisan lapisan ozon adalah pendorong utama
pendinginan stratosfer bawah dari 1979 sampai pertengahan 1990-an.
% ledger: ar6ch3:upper-air

\begin{method}[Sidik jari optimal]\label{met:b2:anthropogenic-climate-change:fingerprinting}
\begin{enumerate}
  \item Simulasikan tanggapan terhadap setiap pendorong secara terpisah: sidik jari $x_1, \dots, x_k$, sebagai
        vektor nilai dalam ruang dan waktu.
  \item Tuliskan pengamatan sebagai $y = \sum_i \beta_i x_i + \varepsilon$, dengan $\varepsilon$
        variabilitas internal, yang statistiknya diambil dari simulasi panjang tanpa gaya paksa.
  \item Perkirakan faktor skala $\beta_i$ dengan kuadrat terkecil, yang diboboti dengan kebalikan
        variabilitas sehingga komponen yang paling berderau paling sedikit diperhitungkan, beserta selang kepercayaannya.
  \item $\beta_i$ yang secara signifikan lebih besar daripada nol berarti bahwa tanggapan terhadap pendorong $i$
        terdeteksi; selang yang memuat 1 dan tidak memuat 0 berarti bahwa tanggapan yang disimulasikan
        konsisten dengan pengamatan, dan perubahan itu dapat diatribusikan sebagian kepada pendorong tersebut.
\end{enumerate}
\end{method}
% ledger: ar6ch3:fingerprint

Bila diterapkan pada suhu permukaan, metode ini memberikan \omterm{def:b2:anthropogenic-climate-change:da}{atribusi} yang diingatkan pada awal bab: dari
\qty{1.06}{\celsius} yang teramati antara 1850--1900 dan 2010--2019, pengaruh manusia \emph{mungkin}
menyebabkan 0.8 sampai \qty{1.3}{\celsius}, dengan perkiraan terbaik \qty{1.07}{\celsius}; pendorong alami dan
variabilitas internal menyumbang tidak lebih dari beberapa persepuluh derajat ke arah mana pun
(\cref{exo:b2:anthropogenic-climate-change:9,exo:b2:anthropogenic-climate-change:10}).
% ledger: ipcc:attribution, ipcc:observed-2010-2019

\section{Skenario dan emisi kumulatif}\label{sec:b2:anthropogenic-climate-change:scenarios}

\begin{definition}[Jalur Sosioekonomi Bersama]\label{def:b2:anthropogenic-climate-change:ssp}
Yang disebut \emph{Jalur Sosioekonomi Bersama}\index{Jalur Sosioekonomi Bersama} (SSP) adalah salah satu dari lima
narasi tentang bagaimana penduduk, perekonomian, dan teknologi mungkin berkembang selama abad ini.
Skenario emisi dalam penilaian itu diberi label SSP$x$-$y$: $x$ menamai jalurnya, dan $y$
gaya radiatif perkiraan, dalam \unit{W.m^{-2}}, yang dicapai skenario itu pada 2100.
\end{definition}
% ledger: ar6spm:ssp

Kelima skenario utama berkisar dari SSP1-1.9 sampai SSP5-8.5. Penilaian sains fisis menerimanya sebagaimana
adanya: penilaian itu tidak menimbang kemungkinannya dan tidak pula menilai narasinya. Untuk masing-masing, model iklim
memberikan proyeksi pemanasan untuk 2081--2100 di atas 1850--1900, dari \qty{1.4}{\celsius} (sangat mungkin 1.0
sampai 1.8) dengan SSP1-1.9 sampai \qty{4.4}{\celsius} (3.3 sampai 5.7) dengan SSP5-8.5. Proyeksi-proyeksi itu mengikuti
satu hukum sederhana yang sama.
% ledger: ar6spm:ssp, ipcc:proj-temp

\begin{definition}[Respons iklim transien terhadap emisi CO$_2$ kumulatif]\label{def:b2:anthropogenic-climate-change:tcre}
Yang disebut \emph{respons iklim transien terhadap emisi \ce{CO2} kumulatif}\index{respons iklim transien terhadap emisi CO$_2$ kumulatif}
(TCRE) adalah pemanasan permukaan global per satuan karbon
dioksida yang diemisikan secara total sejak masa praindustri, dalam \unit{\celsius} per
\qty{1000}{Gt} \ce{CO2}.
\end{definition}

\begin{proposition}[Pemanasan mengikuti emisi kumulatif]\label{prop:b2:anthropogenic-climate-change:tcre}
Dalam pendekatan yang baik, pemanasan yang disebabkan oleh karbon dioksida sebanding dengan jumlah total yang diemisikan,
apa pun lintasan emisinya: $\Delta T \approx \text{TCRE} \times E$, dengan $E$ emisi
kumulatif.
\end{proposition}
\admitted
% ledger: ar6spm:tcre, ar6ch5:tcre-linear

Hukum itu tidak jelas dengan sendirinya. Gaya radiatif hanya tumbuh menurut logaritma konsentrasi, sehingga setiap
ton tambahan menghangatkan lebih sedikit; tetapi samudra menyerap kalor dan karbon makin kurang efisien seiring
emisi yang menumpuk, sehingga lebih banyak bagian setiap ton tetap berada di udara dan pemanasannya kurang tertunda. Kedua
efek itu hampir saling meniadakan, dan semua studi sepakat bahwa itulah asal hukum tersebut. Hubungan itu
hampir linear (\emph{keyakinan tinggi}), dan setiap \qty{1000}{Gt} karbon dioksida \emph{mungkin} menyebabkan
pemanasan 0.27 sampai \qty{0.63}{\celsius}, dengan perkiraan terbaik \qty{0.45}{\celsius}, atau \qty{1.65}{\celsius}
per \qty{1000}{PgC} karbon (\cref{exo:b2:anthropogenic-climate-change:4}). Dari 1850 sampai 2019,
$\num{2390} \pm \qty{240}{Gt}$ karbon dioksida telah diemisikan.
% ledger: ar6ch5:tcre-linear, ar6spm:tcre, ar6spm:budget

\begin{definition}[Sisa anggaran karbon]\label{def:b2:anthropogenic-climate-change:budget}
Yang disebut \emph{sisa anggaran karbon}\index{sisa anggaran karbon} untuk suatu batas suhu adalah
karbon dioksida kumulatif yang masih dapat diemisikan, sejak tanggal yang dinyatakan, dengan pemanasan global tetap
di bawah batas itu pada probabilitas yang dinyatakan, dengan memperhitungkan pemanasan yang disebabkan oleh pendorong-pendorong manusia
lainnya.
\end{definition}

\begin{method}[Sisa anggaran karbon dari TCRE]\label{met:b2:anthropogenic-climate-change:budget}
\begin{enumerate}
  \item Kurangkan pemanasan manusia yang telah tercapai, \qty{1.07}{\celsius} untuk 2010--2019, dari
        batas itu.
  \item Kurangkan pemanasan yang akan ditambahkan pendorong-pendorong lain sampai emisi karbon dioksida
        mencapai nol bersih.
  \item Bagilah dengan TCRE: inilah anggaran sejak pertengahan 2010--2019.
  \item Kurangkan emisi sejak saat itu, dan suatu margin untuk umpan balik sistem Bumi yang lambat, yang
        tidak dimasukkan model (pencairan permafrost adalah yang terbesar).
\end{enumerate}
\end{method}

Penilaian itu menggabungkan lima komponen dengan cara ini: TCRE; pemanasan historis; pemanasan
tambahan oleh pendorong-pendorong lain sekitar 0.1 sampai \qty{0.2}{\celsius}; pemanasan yang mungkin
berlanjut setelah emisi berhenti, yang dinilai nol secara rata-rata (\emph{mungkin} dalam $\pm
\qty{0.19}{\celsius}$); dan umpan balik sistem Bumi, sekitar \qty{26}{Gt} karbon dioksida per
derajat. Untuk \qty{1.5}{\celsius}, yaitu \qty{0.43}{\celsius} di atas 2010--2019, sisa anggaran sejak
awal 2020 adalah \qty{500}{Gt} karbon dioksida pada probabilitas \qty{50}{\percent}, dan
\qty{1350}{Gt} untuk \qty{2}{\celsius}. Pilihan yang dibuat tentang pendorong-pendorong lain menggeser nilai-nilai ini sebesar
\qty{220}{Gt} atau lebih; soal akhir pekan membangun ulang nilai yang pertama.
% ledger: ar6ch5:budget-parts, ar6spm:budget

\begin{dated}{2025}
% ledger: data:co2-mlo, data:gistemp (values read from the datasets as fetched 2026-10-22; offset computed)
Rata-rata tahunan karbon dioksida di Mauna Loa mencapai \qty{427.35}{ppm} pada 2025, suatu gaya radiatif sebesar
$3.93\log_2(427.35/278) = \qty{2.44}{W.m^{-2}}$ menurut hukum logaritmik. Dalam rekaman suhu
NASA, 2024 adalah tahun terhangat, \qty{1.50}{\celsius} di atas rata-rata 1880--1900, tahun-tahun
paling awal rekaman itu; 2025 berada pada \qty{1.40}{\celsius}.
\end{dated}

\begin{omfigure}
\centering
\begin{tikzpicture}
\begin{axis}[omrecord, width=0.86\linewidth, height=5.8cm, xmin=0, xmax=4000, ymin=0, ymax=2.6,
  xlabel={emisi \ce{CO2} kumulatif sejak 1850 (\unit{Gt})},
  ylabel={pemanasan sejak 1850--1900 (\unit{\celsius})}, xtick={0,500,...,4000},
  /pgf/number format/1000 sep={\,}, xtick pos=bottom, ytick pos=left, clip=false,
  tick label style={font=\scriptsize}, label style={font=\small}]
% computed in figdata/bachelor-2/anthropogenic.py part tcre; ledger: ar6spm:tcre, ar6spm:budget
\addplot[name path=hi, draw=none, forget plot] table[x=E_GtCO2, y=dT_hi] {figdata/out/bachelor-2-anthropogenic-tcre.dat};
\addplot[name path=lo, draw=none, forget plot] table[x=E_GtCO2, y=dT_lo] {figdata/out/bachelor-2-anthropogenic-tcre.dat};
\addplot[omThm!18, forget plot] fill between[of=hi and lo];
\addplot[omThm, line width=1.3pt] table[x=E_GtCO2, y=dT_mid] {figdata/out/bachelor-2-anthropogenic-tcre.dat};
\draw[densely dashed, black!60] (axis cs:0,1.5) -- (axis cs:4000,1.5);
\draw[densely dashed, black!60] (axis cs:0,2.0) -- (axis cs:4000,2.0);
\node[anchor=north west, font=\scriptsize, text=black!70] at (axis cs:0,1.5) {\qty{1.5}{\celsius}};
\node[anchor=north west, font=\scriptsize, text=black!70] at (axis cs:0,2.0) {\qty{2}{\celsius}};
\fill[black] (axis cs:2390,1.07) circle (2.2pt);
\draw[black, line width=0.8pt] (axis cs:2150,1.07) -- (axis cs:2630,1.07);
\draw[black, line width=0.8pt] (axis cs:2390,0.8) -- (axis cs:2390,1.3);
\node[anchor=north west, font=\scriptsize, align=left] at (axis cs:2440,1.02) {1850--2019:\\\qty{2390}{Gt}, \qty{1.07}{\celsius}};
\draw[<->, omDef, thick] (axis cs:2390,1.62) -- (axis cs:2890,1.62);
\node[anchor=south, font=\scriptsize, text=omDef, align=center] at (axis cs:2640,1.64) {AR6: \qty{500}{Gt}};
\node[anchor=north west, font=\scriptsize, text=omThm, align=left] at (axis cs:150,2.5) {rentang mungkin TCRE,\\0.27 sampai \qty{0.63}{\celsius} per \qty{1000}{Gt}};
\end{axis}
\end{tikzpicture}
\par\smallskip{\small Pemanasan terhadap emisi karbon dioksida kumulatif untuk TCRE hasil penilaian,
\qty{0.45}{\celsius} per \qty{1000}{Gt} (garis), dan rentang mungkinnya (diarsir), dengan pemanasan manusia
dan emisi sampai saat ini (titik, dengan rentang mungkin sebagai batang). Panah menandai sisa anggaran untuk
\qty{1.5}{\celsius} pada \qty{50}{\percent} sejak 2020: lebih kecil daripada pembacaan garis itu saja,
karena pendorong-pendorong lain juga menghangatkan.}
\end{omfigure}


\section{Muka laut}\label{sec:b2:anthropogenic-climate-change:sea-level}

\begin{definition}[Anggaran muka laut; kenaikan muka laut termosterik]\label{def:b2:anthropogenic-climate-change:sea-level}
Yang disebut \emph{anggaran muka laut}\index{anggaran muka laut} adalah perbandingan antara perubahan muka laut rata-rata
global yang teramati dan jumlah sumbangan-sumbangannya yang diperkirakan: pemuaian samudra yang
menghangat, pencairan gletser serta lembaran es Greenland dan Antarktika, dan perubahan
air yang tersimpan di daratan. Yang disebut \emph{kenaikan muka laut termosterik}\index{kenaikan muka laut termosterik} adalah
bagian yang disebabkan oleh pemuaian termal air laut, pada massa tetap.
\end{definition}

\begin{proposition}[Pemuaian termal suatu lapisan]\label{prop:b2:anthropogenic-climate-change:thermosteric}
Suatu lapisan samudra setebal $H$, dengan koefisien muai termal $\alpha$, yang menghangat secara seragam
sebesar $\Delta T$ pada massa dan luas tetap naik sebesar $\Delta h = \alpha\,\Delta T\,H$.
\end{proposition}
\begin{proof}
Massa $\rho H A$ suatu kolom berluas $A$ konstan, sehingga $\Delta H/H = -\Delta\rho/\rho =
\alpha\,\Delta T$ sampai orde pertama, dengan $\alpha = -(1/\rho)\,\partial\rho/\partial T$.
\end{proof}
% ledger: teos10:alpha-beta

Di samudra yang sebenarnya $\alpha$ bertambah seiring suhu dan tekanan, dan pemanasannya terpusat di dekat
permukaan. Penilaian itu meringkas hasilnya sebagai efisiensi pemuaian sebesar \qty{0.113}{m} per
\qty{e24}{J} kalor yang ditambahkan ke samudra, atau \qty{0.617}{m} per derajat pemanasan seluruh samudra
(\emph{keyakinan sedang}). Atas kedalaman rata-ratanya sebesar \qty{3682}{m}, hal itu berarti $\alpha$ efektif sebesar
$0.617/3682 = \qty{1.68e-4}{\celsius^{-1}}$ (\cref{exo:b2:anthropogenic-climate-change:8}).
% ledger: ar6ch9:expansion, charette:ocean-volume

Anggaran itu, yang digambarkan pada gambar di bawah, tertutup di dalam batas ketidakpastiannya. Selama 1971--2018 sumbangan-sumbangannya berjumlah
\qty{2.00}{mm.yr^{-1}}, dibandingkan dengan kenaikan teramati sebesar \qty{2.33}{mm.yr^{-1}}; selama 2006--2018 jumlahnya
3.61 dibandingkan dengan 3.69. Pemuaian termal adalah suku tunggal terbesar, separuh jumlah selama 1971--2018, tetapi es
mulai menyusul: selama 2006--2018 gletser dan lembaran es bersama-sama memberikan \qty{1.62}{mm.yr^{-1}}, lebih dari
pemuaian itu. Samudra juga terus menghangat di kedalaman lama setelah permukaannya menetap, dan lembaran
es menanggapi selama berabad-abad; muka laut pasti akan naik selama berabad-abad sampai beribu-ribu tahun dan akan
tetap tinggi selama ribuan tahun (\emph{keyakinan tinggi}). Selama dua ribu tahun mendatang muka laut
diproyeksikan naik sekitar 2 sampai \qty{3}{m} jika pemanasan tetap pada \qty{1.5}{\celsius}, 2 sampai
\qty{6}{m} pada \qty{2}{\celsius}, dan 19 sampai \qty{22}{m} pada \qty{5}{\celsius} (\emph{keyakinan rendah}).
% ledger: ar6ch9:table95, ar6spm:sea-commit

\begin{omfigure}
\centering
\begin{tikzpicture}
\begin{axis}[ybar stacked, width=0.72\linewidth, height=6.2cm, ymin=0, ymax=4.3, bar width=26pt,
  xtick=data, symbolic x coords={1971--2018, 1993--2018, 2006--2018}, enlarge x limits=0.25,
  ylabel={laju kenaikan (\unit{mm.yr^{-1}})}, tick label style={font=\scriptsize},
  label style={font=\small}, ytick pos=left, xtick pos=bottom,
  legend style={at={(1.03,0.5)}, anchor=west, font=\scriptsize, draw=black!30},
  legend cell align=left, reverse legend]
% ledger: ar6ch9:table95
\addplot[fill=omDef!70, draw=omDef!90!black] coordinates {(1971--2018,1.01) (1993--2018,1.31) (2006--2018,1.39)};
\addlegendentry{pemuaian termal}
\addplot[fill=omDef!30, draw=omDef!70!black] coordinates {(1971--2018,0.44) (1993--2018,0.55) (2006--2018,0.62)};
\addlegendentry{gletser}
\addplot[fill=omThm!55, draw=omThm!80!black] coordinates {(1971--2018,0.25) (1993--2018,0.43) (2006--2018,0.63)};
\addlegendentry{Greenland}
\addplot[fill=omThm!25, draw=omThm!70!black] coordinates {(1971--2018,0.14) (1993--2018,0.25) (2006--2018,0.37)};
\addlegendentry{Antarktika}
\addplot[fill=black!20, draw=black!50] coordinates {(1971--2018,0.15) (1993--2018,0.31) (2006--2018,0.60)};
\addlegendentry{air daratan}
\node[fill=black, rotate=45, minimum size=5pt, inner sep=0pt] at (axis cs:1971--2018,2.33) {};
\node[fill=black, rotate=45, minimum size=5pt, inner sep=0pt] at (axis cs:1993--2018,3.25) {};
\node[fill=black, rotate=45, minimum size=5pt, inner sep=0pt] at (axis cs:2006--2018,3.69) {};
\end{axis}
\end{tikzpicture}
\par\smallskip{\small \omterm{def:b2:anthropogenic-climate-change:sea-level}{Anggaran muka laut}: sumbangan-sumbangan teramati terhadap kenaikan muka laut rata-rata
global, sebagai laju selama tiga periode yang berakhir pada 2018, ditumpuk, dan laju teramati (belah ketupat). Celah
antara tumpukan dan belah ketupat adalah bagian yang tidak tertutup oleh anggaran itu. Nilai: Laporan Penilaian Keenam
IPCC, Kelompok Kerja~I, Tabel 9.5.}
\end{omfigure}

\section*{Latihan}\label{sec:b2:anthropogenic-climate-change:exercises}

Data untuk latihan, kecuali dinyatakan lain: gaya radiatif penggandaan karbon dioksida $F_{2\times} =
\qty{3.93}{W.m^{-2}}$ dari $C_0 = \qty{278}{ppm}$; tetapan Stefan--Boltzmann $\sigma =
\qty{5.670e-8}{W.m^{-2}.K^{-4}}$; jari-jari Bumi \qty{6371}{km}; satu tahun
\qty{3.156e7}{s}; massa molar karbon dan karbon dioksida 12.011 dan \qty{44.009}{g.mol^{-1}};
umpan balik dan \omterm{def:b2:anthropogenic-climate-change:sea-level}{anggaran muka laut} dalam bab ini.
% ledger: ipcc:erf-2xco2, const:sigma, earth:Rmean, ciaaw:abridged-2024, ar6ch7:table710, ar6ch9:table95

\begin{exercise}[$\star$]\label{exo:b2:anthropogenic-climate-change:1}
Gaya radiatif penggandaan karbon dioksida adalah \qty{3.75}{W.m^{-2}} setelah stratosfer menyesuaikan diri;
troposfer kemudian menyesuaikan diri sebesar $-0.60$ (suhu), +0.22 (uap air), +0.45 (awan), dan
+0.11 (permukaan daratan) \unit{W.m^{-2}}. (a) Tentukan ERF-nya. (b) Dengan persentase berapa nilai itu melebihi
\qty{3.71}{W.m^{-2}} dari penilaian sebelumnya? (c) Penyesuaian mana yang negatif, dan mengapa troposfer yang lebih hangat
mengurangi gaya radiatif?
\end{exercise}
% ledger: ar6ch7:2xco2

\begin{exercise}[$\star$]\label{exo:b2:anthropogenic-climate-change:2}
Sistem iklim memperoleh \qty{434.9}{ZJ} dari 1971 sampai 2018 dan \qty{152.4}{ZJ} dari 2006 sampai 2018,
dihitung di antara pertengahan tahun pertama dan tahun terakhir. (a) Tentukan luas permukaan Bumi.
(b) Konversikan setiap perolehan menjadi ketidakseimbangan rata-rata dalam \unit{W.m^{-2}}. (c) Samudra mengambil
\qty{91}{\percent}: berapa zettajoule selama 1971--2018?
\end{exercise}
% ledger: ar6ch7:inventory, earth:Rmean

\begin{exercise}[$\star$]\label{exo:b2:anthropogenic-climate-change:3}
(a) Jumlahkan nilai tengah umpan balik pada gambar, dan bandingkan dengan nilai bersihnya. (b) Tentukan
ECS dengan nilai bersih itu, dan pemanasan dengan \omterm{def:b2:anthropogenic-climate-change:feedbacks}{respons Planck} saja. (c) Dengan faktor berapa
umpan balik lainnya memperkuat pemanasan Planck?
\end{exercise}

\begin{exercise}[$\star$]\label{exo:b2:anthropogenic-climate-change:4}
(a) Konversikan TCRE sebesar \qty{0.45}{\celsius} per \qty{1000}{Gt} karbon dioksida menjadi derajat per
\qty{1000}{PgC} karbon, dan lakukan hal yang sama untuk rentang mungkinnya, 0.27 sampai 0.63. (b) Konversikan
\qty{2390}{Gt} karbon dioksida yang diemisikan dari 1850 sampai 2019 menjadi petagram karbon. (c) Berapa
pemanasan yang diberikan TCRE untuk emisi ini?
\end{exercise}
% ledger: ar6spm:tcre, ar6spm:budget

\begin{exercise}[$\star\star$]\label{exo:b2:anthropogenic-climate-change:5}
(a) Permukaan memancarkan \qty{398}{W.m^{-2}} dan efek rumah kaca menahan \qty{159}{W.m^{-2}} darinya.
Tentukan $g$ dan suhu permukaan benda hitam yang memancarkan \qty{398}{W.m^{-2}}. (b) Tentukan
\omterm{def:b2:anthropogenic-climate-change:feedbacks}{respons Planck} pada \qty{288}{K}, dan pada \qty{298}{K} dengan $g$ yang sama. (c) Apakah \omterm{def:b2:anthropogenic-climate-change:feedbacks}{respons
Planck} saja akan membuat Bumi yang lebih hangat lebih peka atau kurang peka?
\end{exercise}
% ledger: ipcc:budget-surface, ar6ch7:planck

\begin{exercise}[$\star\star$]\label{exo:b2:anthropogenic-climate-change:6}
Satelit memberikan \omterm{def:b2:anthropogenic-climate-change:feedbacks}{umpan balik uap air} sebesar $1.85 \pm 0.32$ dan studi-studi memberikan \omterm{def:b2:anthropogenic-climate-change:feedbacks}{umpan balik laju lapse} sebesar
$-0.50 \pm 0.20$ \unit{W.m^{-2}.\celsius^{-1}} (satu simpangan baku). (a) Tentukan jumlahnya.
(b) Jika keduanya saling bebas, berapa simpangan baku jumlah itu? (c) Menurut model-model,
simpangan baku jumlah itu hanya \num{0.10}. Apa yang dikatakan hal itu tentang kedua galat tersebut?
\end{exercise}
% ledger: ar6ch7:wv-lr

\begin{exercise}[$\star\star$]\label{exo:b2:anthropogenic-climate-change:7}
Dari 1850--1900 sampai 2006--2019, $\Delta N = \qty{0.59}{W.m^{-2}}$, $\Delta T = \qty{1.03}{\celsius}$,
dan $\Delta F = \qty{2.20}{W.m^{-2}}$. (a) Tentukan $\lambda$ dan sensitivitas efektifnya. (b) Parameter
umpan balik diperkirakan menjadi kurang negatif sebesar $\lambda' = \qty{0.5}{W.m^{-2}.\celsius^{-1}}$
seiring berkembangnya pola pemanasan. Tentukan sensitivitas $-F_{2\times}/(\lambda + \lambda')$.
(c) Ulangi (b) dengan $\lambda' = 0.1$.
\end{exercise}
% ledger: ar6ch7:energy-budget, ar6ch7:pattern, ar6ch7:pattern-ecs

\begin{exercise}[$\star\star$]\label{exo:b2:anthropogenic-climate-change:8}
(a) Dari kenaikan \qty{0.617}{m} per derajat pemanasan seluruh samudra, dengan kedalaman rata-rata
\qty{3682}{m}, tentukan koefisien muai efektifnya. (b) Samudra memperoleh \qty{396}{ZJ} selama
1971--2018. Dengan efisiensi pemuaian \qty{0.113}{m} per \qty{e24}{J}, tentukan kenaikan
termosteriknya, dan bandingkan dengan \qty{47.5}{mm} dalam anggaran itu. (c) Suatu lapisan setebal \qty{700}{m} menghangat sebesar
\qty{0.2}{\celsius} dengan $\alpha = \qty{2e-4}{\celsius^{-1}}$: berapa besar lapisan itu menaikkan muka laut?
\end{exercise}
% ledger: ar6ch9:expansion, charette:ocean-volume, ar6ch7:inventory, ar6ch9:table95

\begin{exercise}[$\star\star$]\label{exo:b2:anthropogenic-climate-change:9}
Suatu sidik jari model, dalam empat dekade, adalah $x = (0.10, 0.20, 0.30, 0.40)$ \unit{\celsius};
pengamatannya adalah $y = (0.12, 0.18, 0.33, 0.41)$ \unit{\celsius}. (a) Tentukan faktor skala kuadrat terkecil
$\beta$ dari $y = \beta x + \varepsilon$. (b) Galat bakunya, yang diperkirakan dari simulasi tanpa gaya paksa,
adalah \num{0.08}. Apakah tanggapan itu terdeteksi? Apakah tanggapan itu konsisten dengan model? (c) Apa
arti $\beta$ sebesar 0.5 dengan galat yang sama?
\end{exercise}

\begin{exercise}[$\star\star\star$]\label{exo:b2:anthropogenic-climate-change:10}
Dua sidik jari, gas rumah kaca $g = (0.2, 0.4, 0.6, 0.8)$ dan aerosol $a = (-0.1, -0.2,
-0.2, -0.3)$, keduanya dalam \unit{\celsius} selama empat dekade, dicocokkan dengan pengamatan $y = (0.1,
0.2, 0.4, 0.5)$. (a) Tuliskan persamaan normal untuk pencocokan kuadrat terkecil $y = \beta_g g + \beta_a a
+ \varepsilon$, lalu selesaikan. (b) Tentukan pemanasan yang diatribusikan kepada setiap pendorong pada dekade terakhir.
(c) Mengapa ketidakpastian $\beta_g$ dan $\beta_a$ membesar bila kedua sidik jari itu tampak mirip?
\end{exercise}

\begin{exercise}[$\star\star\star$]\label{exo:b2:anthropogenic-climate-change:11}
Gunakan metode \omterm{def:b2:anthropogenic-climate-change:budget}{sisa anggaran karbon} dengan TCRE sebesar \qty{0.45}{\celsius} per
\qty{1000}{Gt}, \qty{210}{Gt} yang diemisikan dari 2015 sampai 2019, dan \qty{26}{Gt} per derajat untuk
umpan balik sistem Bumi. (a) Tentukan anggaran untuk \qty{2}{\celsius}, yaitu \qty{0.93}{\celsius} di atas
2010--2019, jika pendorong-pendorong lain menambahkan \qty{0.1}{\celsius}. (b) Pemanasan tambahan berapa oleh pendorong-pendorong
lain yang membuatnya sama dengan \qty{1350}{Gt} hasil penilaian? Apakah nilai itu di dalam rentang hasil penilaian? (c) Tunjukkan
bahwa setiap \qty{0.1}{\celsius} pemanasan non-CO$_2$ mengubah anggaran mana pun dengan besar yang sama.
\end{exercise}

\begin{exercise}[$\star\star\star$]\label{exo:b2:anthropogenic-climate-change:12}
Umpan balik bersih memiliki rentang sangat mungkin dari $-1.81$ sampai \qty{-0.51}{W.m^{-2}.\celsius^{-1}}, di sekitar
$-1.16$. (a) Tentukan ECS pada kedua ujungnya dan di tengahnya. (b) Tunjukkan bahwa rentang $\lambda$ yang
simetris terhadap tengahnya memberikan rentang ECS yang tidak simetris, dan sebutkan sisi mana yang lebih panjang. (c) Mengapa
penilaian-penilaian itu lebih sulit menyingkirkan ECS yang tinggi daripada ECS yang rendah?
\end{exercise}
% ledger: ar6ch7:table710

\begin{problem}[{Soal akhir pekan --- Membaca anggaran}]\label{pb:b2:anthropogenic-climate-change:1}
Soal ini mengikuti energi dan karbon dalam penilaian itu dari gaya radiatif sampai sisa
anggaran untuk \qty{1.5}{\celsius}. Gunakan data latihan, dan: ERF 1750--2019 (\unit{W.m^{-2}})
karbon dioksida 2.16, metana 0.54, dinitrogen oksida 0.21, gas terhalogenasi 0.41, ozon 0.47,
uap air stratosfer 0.05, aerosol--radiasi $-0.22$, aerosol--awan $-0.84$, tata guna lahan
$-0.20$, karbon hitam di atas salju 0.08, jejak kondensasi 0.06; karbon dioksida \qty{410}{ppm} pada 2019 dan
\qty{427.35}{ppm} pada 2025; pemanasan \qty{0.79}{W.m^{-2}} selama 2006--2018; perubahan anggaran energi
$\Delta N = 0.59$, $\Delta F = \qty{2.20}{W.m^{-2}}$, $\Delta T = \qty{1.03}{\celsius}$; TCRE
\qty{0.45}{\celsius} per \qty{1000}{Gt} karbon dioksida (mungkin 0.27 sampai 0.63); \qty{2390}{Gt}
diemisikan 1850--2019; pemanasan manusia \qty{1.07}{\celsius} untuk 2010--2019; \qty{210}{Gt} diemisikan dari
2015 sampai 2019; pemanasan non-CO$_2$ tambahan \qty{0.1}{\celsius}, ujung bawah rentang hasil penilaian 0.1
sampai \qty{0.2}{\celsius}; umpan balik sistem Bumi \qty{26}{Gt} per derajat. Bulatkan sampai tiga angka
penting.
% ledger: ipcc:erf-agents, ipcc:erf-2xco2, data:co2-mlo, ar6ch7:inventory, ar6ch7:energy-budget, ar6spm:tcre, ar6spm:budget, ar6ch5:budget-parts

\noindent\textbf{Bagian I --- Gaya radiatif.}
\begin{enumerate}
  \item Jumlahkan gaya radiatif keempat gas rumah kaca yang tercampur baik.
  \item Jumlahkan gaya radiatif aerosol.
  \item Jumlahkan semua suku manusia. Bandingkan dengan \qty{2.72}{W.m^{-2}} hasil penilaian.
  \item Periksalah gaya radiatif karbon dioksida pada 2019 dengan hukum logaritmik.
  \item Tentukan gaya radiatif karbon dioksida pada 2025 dengan hukum yang sama.
  \item Suku mana yang ketidakpastiannya paling lebar, dan mengapa suku itu mendominasi ketidakpastian total?
\end{enumerate}

\noindent\textbf{Bagian II --- Anggaran energi.}
\begin{enumerate}[resume]
  \item Berapa banyak energi yang diperoleh sistem iklim selama 2006--2018, dalam zettajoule?
  \item Berapa banyak darinya yang masuk ke samudra?
  \item Tentukan parameter umpan balik dari perubahan anggaran energi.
  \item Tentukan \omterm{def:b2:anthropogenic-climate-change:ecs}{sensitivitas iklim kesetimbangan} efektif.
  \item Tentukan \omterm{def:b2:anthropogenic-climate-change:tcr}{respons iklim transien} dari data yang sama.
  \item Tentukan ECS dari umpan balik bersih pada gambar, dan bandingkan kedua nilai itu.
  \item Berikan alasan fisis mengapa perkiraan historis adalah yang lebih rendah.
\end{enumerate}

\noindent\textbf{Bagian III --- Pemanasan per ton.}
\begin{enumerate}[resume]
  \item Konversikan TCRE menjadi derajat per \qty{1000}{PgC} karbon.
  \item Konversikan rentang mungkinnya.
  \item Konversikan emisi 1850--2019 menjadi petagram karbon.
  \item Berapa pemanasan yang diberikan TCRE untuk emisi 1850--2019?
  \item Bandingkan dengan pemanasan manusia sebesar \qty{1.07}{\celsius}. Apa yang dikatakan perbandingan itu tentang
        efek bersih pendorong-pendorong manusia lainnya?
\end{enumerate}

\noindent\textbf{Bagian IV --- Sisa anggaran untuk \qty{1.5}{\celsius}.}
\begin{enumerate}[resume]
  \item Berapa pemanasan yang tersisa di atas 2010--2019?
  \item Kurangkan pemanasan non-CO$_2$ dan bagilah dengan TCRE: berapa anggaran yang tersisa sejak pertengahan
        2010--2019?
  \item Kurangkan emisi 2015--2019.
  \item Kurangkan umpan balik sistem Bumi atas sisa pemanasan.
  \item Seberapa besar anggaran itu berubah jika pendorong-pendorong lain menambahkan \qty{0.2}{\celsius} sebagai gantinya?
  \item Sebagai ringkasan: berapa \omterm{def:b2:anthropogenic-climate-change:budget}{sisa anggaran karbon} untuk \qty{1.5}{\celsius}, sejak awal
        2020, yang dihitung dari TCRE? Bandingkan dengan \qty{500}{Gt} hasil penilaian.
\end{enumerate}
\end{problem}
